\chapter{Aprendizaje No Supervisado: Segmentación de Mercado}

\section{Metodología de Clustering}
Antes de acometer la tarea predictiva, resultaba imperativo comprender si la oferta de alquileres en Nueva York obedecía a distribuciones homogéneas o si, por el contrario, existían submercados latentes. Se empleó el algoritmo \textbf{K-Means} operando estrictamente sobre el espacio latente reducido (generado por PCA), garantizando así que las distancias Euclidianas calculadas por el algoritmo estuviesen exentas del sesgo impuesto por magnitudes dispares o colinealidad intrínseca.

La función de coste a minimizar (Inercia o WCSS - \textit{Within-Cluster Sum of Squares}) se define como:
\begin{equation}
J = \sum_{j=1}^{k} \sum_{i=1}^{n} ||x_i^{(j)} - c_j||^2
\end{equation}

\section{Determinación por Consenso del Número Óptimo de Clústeres (K)}
La dependencia analítica a un único método heurístico (e.g., Regla del Codo) para la selección de $K$ puede inducir sesgos visuales e interpretativos severos. Este proyecto apostó por una metodología de validación tripartita y unificada:
\begin{enumerate}
    \item \textbf{Método del Codo (Elbow Method):} Análisis del punto de máxima curvatura sobre la curva de inercia decreciente, identificando el punto de saturación donde la ganancia marginal de añadir centroides adicionales colapsa.
    \item \textbf{Análisis de Silueta (Silhouette Score):} Maximización de la cohesión intra-cluster frente a la separación inter-cluster.
    \item \textbf{Estadístico Gap (Gap Statistic):} Comparación de la dispersión de los datos frente a una distribución nula aleatoria de referencia.
\end{enumerate}
A partir de estos tres cálculos, el algoritmo extrajo dinámicamente la **mediana estadística** (\textit{Consensus K}) de las tres predicciones heurísticas, fijando la partición de la base de datos de manera agnóstica a intervenciones manuales.

\section{Perfilado y Caracterización}
La segmentación resultante arrojó perfiles de mercado sustancialmente disímiles y dotados de semántica interpretable. La inyección de inteligencia analítica sobre los resultados permitió aislar las modas estadísticas de la variable \texttt{room\_type} por cada clúster. 

Los hallazgos demuestran que el mercado neoyorquino está fuertemente fragmentado por una conjunción bivariante entre el tipo de habitación ofertada y el eje centralidad-precio. Se identificaron clústeres representativos de **ofertas premium** (habitaciones de tipo \textit{Entire home/apt} de alto costo ubicados a mínimas distancias de los hitos espaciales) frente a **clústeres residuales** (mayoritariamente dominados en porcentaje por opciones \textit{Shared room} o \textit{Private room} en las periferias del Bronx o zonas externas de Queens y Brooklyn). 

Este hallazgo es crucial: dictamina que una estrategia estandarizada de precios (\textit{One-size-fits-all}) fracasaría inevitablemente, justificando la necesidad subyacente de emplear algoritmos de regresión flexibles capaces de emular ramificaciones condicionales.
